\subsection{The exact weight-nine normalization}
\label{subsec:weight-nine-exact}

The integral boundary measure forces the ninth coefficient into
its arithmetic lattice.

\begin{theorem}\label{thm:weight-nine-exact-image}
One has
\[
 m_9(D_8)=m_9(J_8)=512\mathbf Z_2.
\]
The same equality holds with $D_8,J_8$ replaced by $D_7,J_7$.
For $D_9=D_8\cap\ker m_9$ and $J_9=D_9\cap H$,
\[
 D_8=D_9J_8,\qquad G=D_9H,\qquad
 B\simeq D_9/\overline{[G,D_9]J_9}.
\]
The subgroup $D_9$ is proper, closed and normal in $G$, and every
word belonging to it starts in degree at least ten. Moreover,
\[
 (m_8,m_9):D_7\longrightarrow\mathbf Z_2\times512\mathbf Z_2
\]
is surjective already on $J_7$ and has kernel $D_9$.
\end{theorem}

\begin{proof}
The arithmetic word in Proposition~\ref{prop:weight-nine-arithmetic-value}
gives a value of exact valuation nine on $J_8$.
For the upper bound, use Lemma~\ref{lem:uniform-odd-coordinate-bound}
with $r=4$. On $D_7$ the coefficients $a_2,a_4,a_6$ vanish, so
$m_9=a_8\in512\mathbf Z_2$. Concretely, the polynomial
\[
 (t^2-1)(t^2-9)(t^2-25)(t^2-49)
 =t^8-84t^6+1974t^4-12916t^2+11025
\]
is divisible by $2^{15}$ on odd $t$. Its lower moments vanish,
and its eighth moment is $(1-2^8)\tau_9$, with
$\tau_9=20160m_9$ and $v_2(20160)=6$. This is exactly the required
valuation-nine bound. Compactness and the arithmetic witness
therefore prove equality of the two images.

The coefficient $m_9=c_9$ is already an additive, weight-nine
homomorphism on $D_7$. Given any element of $D_7$, multiply it by
a suitable profinite power of the primitive arithmetic
weight-eight commutator $w_8$ to make $m_8$ zero. This does not
change $c_9$, because every logarithmic beta homomorphism vanishes
on $w_8$. Consequently the same exact image equality holds on
$D_7$ and $J_7$. The two arithmetic words $w_8,w_9$ have respective
coordinate values
\[
 (m_8(w_8),0),\qquad(0,m_9(w_9)),
\]
where the first is a unit and the second has exact valuation
nine. They prove the asserted joint surjectivity.

The scalar conjugation laws imply normality of $D_9$. The
arithmetic value proves properness, and equality of the two images
gives $D_8=D_9J_8$, hence $G=D_9H$. The same compact gluing and
centralization argument as in the preceding normalizations gives
the displayed description of $B$. The leading-line calculation
shows that $m_9=0$ kills the whole degree-nine term, so the words
of $D_9$ start in degree at least ten.
\end{proof}

\subsection{The ninth character and its extension obstruction}

\begin{proposition}\label{prop:eta-nine-nonextension}
The character
\[
 \eta_9:D_7\longrightarrow\mathbf F_2,\qquad
 \eta_9(g)=m_9(g)/512\pmod2
\]
is continuous and $G$-invariant. Its restriction to $D_8$ has
nonzero arithmetic restriction and has no $G$-invariant extension
to $D_5$.
\end{proposition}

\begin{proof}
Integrality on $D_7$ and invariance follow from the exact image
and the weight-nine law. Recall the arithmetic construction
\[
 z_7=h z_5h^{-1}z_5^{-t^7}\in J_7,
 \qquad z_5\in J_5,
 \qquad w_9=z_7w_8^{-u}\in J_8,
\]
where $u_2=m_8(z_7)/m_8(w_8)$ and $m_8(w_8)$ is a unit.
The invariant character $\alpha_8$ on $D_5$ vanishes on $z_7$,
since $t$ is odd. Its restriction to $D_7$ is $m_8\bmod2$, so
$u_2$ is even.

Let $\chi:D_5\to\mathbf F_2$ be $G$-invariant. The same odd-power
argument gives $\chi(z_7)=0$. Also
$\chi(w_8^{-u})=-u_2\chi(w_8)=0$. Hence $\chi(w_9)=0$, whereas
$\eta_9(w_9)=1$ by its exact arithmetic valuation. This rules out
an invariant extension from $D_8$ to $D_5$.
\end{proof}

\begin{corollary}\label{cor:weight-nine-character-kernels}
The invariant quadratic-character kernels on restriction to $D_9$
are as follows:
\[
\begin{array}{c|l}
 \text{domain}&\text{kernel}\\ \hline
 D_8&\mathbf F_2\eta_9\\
 D_7&\mathbf F_2\eta_8\oplus\mathbf F_2\eta_9\\
 D_5&\mathbf F_2\eta_7\oplus\mathbf F_2\alpha_8\\
 D_3&\mathbf F_2\eta_5\\
 N&\displaystyle\left(\bigoplus_{q\text{ odd prime}}\mathbf F_2\psi_q\right)
          \oplus\mathbf F_2\eta_3.
\end{array}
\]
\end{corollary}

\begin{proof}
The first two rows follow from the exact scalar and joint quotient
maps. A character on $D_5$ killed on $D_9$ restricts on $D_8$ to
zero or $\eta_9$; the latter is impossible by the preceding
proposition. Its restriction therefore already kills $D_8$,
whose kernel on $D_5$ is the displayed span of $\eta_7,\alpha_8$.
No nonzero combination of those two characters extends to $D_3$,
so the $D_3$ kernel is unchanged. Similarly any invariant character
on $N$ killed on $D_9$ must kill $D_8$, since otherwise its
restriction to $D_5$ would extend $\eta_9$. The preceding kernel
formula for $N$ then gives the last row.
\end{proof}
